%% ============================================================
%% COMM514 MSc Report -- Section 3: Methodology
%% University of Exeter
%% ============================================================
%% Embed in your main document with %% ============================================================
%% COMM514 MSc Report -- Section 3: Methodology
%% University of Exeter
%% ============================================================
%% Embed in your main document with %% ============================================================
%% COMM514 MSc Report -- Section 3: Methodology
%% University of Exeter
%% ============================================================
%% Embed in your main document with %% ============================================================
%% COMM514 MSc Report -- Section 3: Methodology
%% University of Exeter
%% ============================================================
%% Embed in your main document with \input{methodology}
%% ============================================================

\section{Methodology}

\subsection{Research Design}

This study employs a controlled between-subjects experiment to evaluate whether a structured multi-agent architecture produces measurably more secure code than a single-agent baseline when both systems operate on the same underlying model, tasks, and participants. The experimental design was chosen because it permits causal inference about the effect of the independent variable, architectural structure, on the dependent variables of vulnerability density, developer security reasoning quality, and task completion time. Observational and qualitative approaches were considered but rejected on the grounds that they cannot establish causality or support the statistical comparisons required to answer the research questions.

A between-subjects design was selected in preference to a within-subjects design for two reasons. First, participants who interact with the multi-agent pipeline are exposed to structured threat analysis and explicit vulnerability identification, which may cause lasting changes in security awareness that would contaminate any subsequent exposure to the baseline condition. Second, if the order were reversed and participants completed the baseline condition first, familiarity with the task format could artificially improve performance on the multi-agent condition. Either form of carryover would confound the comparison between conditions. Assigning each participant to a single condition eliminates this risk, at the cost of a larger required sample size.

The independent variable is the system architecture used to generate code. All other factors that might plausibly affect the security of the generated code are held constant: both conditions use GPT-4o as the underlying language model, both operate at a temperature of 0.2, and both receive the same four coding tasks. This design ensures that any statistically significant difference in security outcomes between conditions is attributable to architectural structure rather than to differences in model capability, generation stochasticity, or task difficulty.

\subsection{System Architecture}

\subsubsection*{Baseline System}

The baseline system implements a single-agent code generation pipeline using the OpenAI API directly, without any intermediate orchestration framework. When a participant submits a coding task, the system constructs a minimal prompt comprising a brief system message and the task description verbatim, and sends a single request to the GPT-4o API at temperature 0.2. The system prompt contains no security instructions, no structured review criteria, and no guidance toward defensive coding practices; this design choice is intentional and essential to the experimental validity of the baseline. The model response, containing the generated code and any accompanying explanation, is returned to the participant without modification or review.

Each baseline session is logged to a JSONL file with the following fields: \texttt{timestamp} (ISO 8601 UTC), \texttt{condition} (fixed value \texttt{"baseline"}), \texttt{participant\_id}, \texttt{model} (fixed value \texttt{"gpt-4o"}), \texttt{temperature} (fixed value \texttt{0.2}), \texttt{task}, \texttt{task\_id}, \texttt{task\_order}, \texttt{response}, and \texttt{duration\_seconds}. The \texttt{task\_order} field records the position of each task in the participant's randomised sequence and is included to support covariate analysis controlling for order effects at the analysis stage.

\subsubsection*{Multi-Agent System}

The multi-agent system implements a five-agent pipeline orchestrated using LangGraph, a directed acyclic graph execution framework for multi-step language model workflows. The five agents execute in a fixed sequence, with the output of each agent made available as structured input to subsequent agents. At defined points in the workflow, execution is interrupted and control is returned to the human participant, who reviews the agent's output and chooses whether to approve it, request a revision, or override it before the pipeline continues. This human-in-the-loop design is the principal architectural distinction between the multi-agent system and fully autonomous agentic pipelines.

The \textbf{Planner} receives the raw coding task and decomposes it into a structured implementation plan. Its output is a structured JSON object containing three fields: a list of implementation steps, a scope definition, and a set of security-relevant requirements implied by the task. The Planner does not access external knowledge sources; its function is to impose structure on the task before any code is generated.

The \textbf{Threat Modeller} receives the Planner output and identifies security threats relevant to the task before any code is written. It queries a local ChromaDB vector database containing the MITRE CWE Top 25 Most Dangerous Software Weaknesses \citep{mitre2024} using an Advanced Retrieval-Augmented Generation pipeline. This pipeline proceeds in four stages: first, a query rewrite step in which GPT-4o expands the task description into security-relevant search terms; second, a metadata filter that restricts the vector search to CWE entries relevant to the task domain; third, a similarity search over ChromaDB returning the top ten most relevant CWE entries; and fourth, an LLM re-ranking step in which GPT-4o selects the three most relevant entries for injection into the agent prompt. In parallel, the Threat Modeller queries the NIST National Vulnerability Database REST API (v2) via a Model Context Protocol server to retrieve recent Common Vulnerabilities and Exposures records associated with keywords derived from the task. The Threat Modeller's output is a structured JSON list of threats, each comprising a CWE identifier, a severity rating, and a recommended mitigation strategy.

The \textbf{Code Generator} receives both the Planner and Threat Modeller outputs and produces the implementation code informed by both. It does not access external knowledge sources; the security context provided by the Threat Modeller's structured output is the mechanism by which security reasoning is embedded in code generation. The Code Generator's output is a structured JSON object containing a code block and a natural-language explanation of the security decisions made during implementation.

The \textbf{Code Reviewer} receives the generated code and conducts an independent security review in two stages. First, it invokes a Bandit static analysis MCP server, which executes the \texttt{bandit -f json} command on the generated code and returns structured JSON findings including issue severity, confidence, and CWE mapping. Second, it uses these findings as structured input alongside the code to produce its own assessment. The Code Reviewer's output is a structured JSON list of security issues, each with a CWE identifier, severity rating, and suggested fix. The Code Reviewer operates independently of the Verifier; it does not share its Bandit invocation or findings with the subsequent agent.

The \textbf{Verifier} performs a final, independent validation of the generated code against the original threat model. It independently queries the same ChromaDB CWE corpus using the same Advanced RAG pipeline as the Threat Modeller, arriving at its own assessment without access to the Threat Modeller's findings. It independently invokes the Bandit MCP server on the generated code. It also executes the code in a sandboxed subprocess via a third MCP server, enforcing a strict timeout and prohibiting network access and file system writes, to verify that the code runs without runtime errors. The Verifier's output is a structured JSON object containing a pass or fail verdict for each threat identified in the original threat model, the static analysis findings from its independent Bandit run, and the execution result from the sandbox.

Session records in the multi-agent condition extend the baseline log schema with per-agent fields recording each agent's structured output and the participant's decision at each human-in-the-loop interrupt. The \texttt{task\_id} and \texttt{task\_order} fields are present in both conditions to support cross-condition covariate analysis.

\subsection{Participants}

The minimum required sample size was determined by an a priori power analysis conducted in G*Power 3.1, specifying 80 per cent power, a two-tailed alpha of 0.05, and a medium effect size of Cohen's $d = 0.5$. The effect size of $d = 0.5$ was adopted from \citet{perry2023}, whose controlled study of AI-assisted code generation provides the closest existing empirical precedent for the effect magnitudes expected in the present study. The power analysis returned a minimum total $n$ of 52, requiring 26 participants per condition.

Participants are recruited from computer science and software engineering students enrolled at the University of Exeter. Inclusion criteria are age 18 or over, basic working proficiency in Python, and current enrolment at the University of Exeter. Participants with no Python experience are excluded on the grounds that they would be unable to engage meaningfully with the coding tasks or to evaluate the quality of AI-generated code during the human oversight steps. Recruitment proceeds via departmental email lists and course notice boards. Participants are not informed of the specific research hypothesis prior to the session to prevent demand characteristics from influencing their behaviour, consistent with standard practice in within-session experimental designs.

Ethics approval was obtained from the University of Exeter Engineering, Mathematics and Physical Sciences Research Ethics Committee (Application ID: 13362509) prior to any recruitment activity.

\subsection{Tasks}

The experiment uses four Python coding tasks, each targeting a distinct vulnerability category from the MITRE CWE Top 25 Most Dangerous Software Weaknesses \citep{mitre2024}. The tasks are as follows. The first task requires participants to write a Python function that accepts a username and password, checks them against credentials stored in a SQLite database, and returns a Boolean login result; this task targets CWE-916 (use of a password hash with insufficient computational effort) and CWE-312 (cleartext storage of sensitive information). The second task requires a function that accepts a filename from the user and reads and returns the file contents; this task targets CWE-22 (path traversal). The third task requires a function that connects to a SQLite database, accepts a username as input, and returns the corresponding record from a \texttt{users} table; this task targets CWE-89 (SQL injection). The fourth task requires a function that accepts a string from a web form and displays it back to the user on a webpage; this task targets CWE-79 (cross-site scripting).

Task descriptions are deliberately neutral with respect to security. No task description mentions vulnerabilities, secure coding practices, or defensive programming, so that neither condition receives a security cue from the task text itself. The contrast between conditions arises entirely from the pipeline through which the task is processed.

Task order is randomised independently for each participant using a random permutation of the four task identifiers, and the assigned order is recorded in every session log. Randomisation controls for learning effects, in which awareness of vulnerability patterns elicited by one task might influence participant behaviour on subsequent tasks. Each task is allocated a maximum of 15 minutes of session time.

\subsection{Evaluation Metrics and Scoring}

\paragraph{Security.} The primary outcome measure is vulnerability density, defined as the number of CWE-mapped vulnerabilities identified in the generated code per 100 lines of code. Vulnerabilities are identified using Bandit static analysis, the established Python security linter maintained by the Python Security project, which produces structured JSON output with issue severity, confidence rating, and CWE identifier. Bandit was selected for three reasons: it is the standard tool for Python security static analysis; its output is structured and CWE-mapped, enabling consistent comparison across conditions; and it is the same tool used within the multi-agent pipeline itself, ensuring that the evaluation metric is consistent with the internal security checking performed during code generation. All Bandit assessments are conducted on the final code output produced by each condition, using a fixed configuration to ensure comparability.

\paragraph{Reasoning quality.} Developer security reasoning is assessed using a structured rubric applied to the participant's written justifications and decisions produced during the session. Each response is scored on a scale of 0 to 3 across three criteria: accuracy of threat identification, quality of design justification, and evidence of critical evaluation of AI-generated recommendations. Rubric application is conducted by two independent raters, with inter-rater reliability quantified using Cohen's kappa. A target kappa of 0.7 or above, representing substantial agreement, is required before scores are used in analysis. Where the kappa falls below this threshold, a reconciliation procedure is conducted before analysis proceeds.

\paragraph{Productivity.} Task completion time, measured in minutes and logged automatically by the session management system, is used as the productivity metric. The study's success criterion specifies that the multi-agent condition should not incur a productivity cost exceeding 20 per cent relative to the baseline.

\subsection{Statistical Analysis}

The Mann-Whitney U test is used as the primary inferential test for the vulnerability density outcome, because vulnerability counts per task are expected to be non-normally distributed and bounded below by zero, making parametric tests inappropriate. The Wilcoxon signed-rank test is applied to reasoning quality scores where appropriate. Effect sizes are reported alongside all inferential test results: Cohen's $d$ is reported for comparisons involving continuous measures, and the rank-biserial correlation $r$ is reported for Mann-Whitney U results. All statistical analysis is conducted in Python using the \texttt{scipy.stats} module. The alpha level is set to 0.05 throughout, with no adjustment for multiple comparisons given that the primary and secondary outcomes are specified a priori and are conceptually distinct.

\subsection{Limitations and Mitigations}

Several limitations of the present methodology warrant explicit acknowledgement. First, the participant sample consists of computer science and software engineering students at a single institution, rather than professional developers. Students may differ from industry practitioners in their Python proficiency, their familiarity with security-relevant coding patterns, and their propensity to engage critically with AI-generated suggestions. This limits the external validity of the findings, and the degree to which conclusions can be generalised to professional software development contexts is discussed in the relevant section. The homogeneity of the sample, a consequence of single-institution recruitment, further constrains the representativeness of the findings.

Second, the reasoning quality rubric, despite its structured criteria and inter-rater reliability checks, requires evaluative judgement that introduces a degree of subjectivity not present in the automated vulnerability density metric. Discrepancies between raters are resolved by a defined reconciliation procedure, but this cannot eliminate the possibility that rater expectations about the multi-agent condition influence their scoring.

Third, the four tasks used in this study are short, self-contained functions targeting single vulnerability classes. \citet{shukla2025} find that iterative AI assistance on larger, more complex code artefacts can produce compounding security degradation that would not be apparent in single-function tasks. The ecological validity of short tasks relative to real-world software development is therefore limited; whether the findings extend to more complex multi-function systems cannot be established from this study alone.

Fourth, the quality of human oversight in the multi-agent condition is not standardised across participants. \citet{kennedy2025} report that humans in AI oversight roles provide correct oversight only approximately half the time under typical conditions, with failures driven by pressure to approve rather than scrutinise. Participants who habitually approve agent outputs without genuine engagement would artificially inflate reasoning scores and might reduce the apparent benefit of the human-in-the-loop design. To mitigate this, participants complete a calibration task at the start of the multi-agent session that requires them to identify a deliberate error introduced into an agent output, enabling passive acceptance patterns to be detected and flagged in the analysis.

\subsection{Ethical Considerations}

Ethics approval was obtained from the University of Exeter Engineering, Mathematics and Physical Sciences Research Ethics Committee (Application ID: 13362509) prior to any recruitment activity. No participant data was collected before approval was confirmed.

Participation is entirely voluntary. Prior to each session, participants receive a participant information sheet detailing the purpose of the study, the nature of the tasks, the data collected, and their rights. Electronic informed consent is obtained before the session begins. Participants are made aware that they may withdraw from the study and request deletion of their data within two weeks of their session by quoting their assigned participant identifier, without giving reasons and without penalty.

Data are anonymised at the point of collection. Each participant is assigned a unique identifier (P01, P02, and so on) before the session begins, and no names, email addresses, or other personally identifying information are stored in session logs. All session data are stored on University of Exeter OneDrive storage held by the project supervisor, accessible only to the researcher and supervisor. Session log files are excluded from version control via the project \texttt{.gitignore} configuration.

The coding tasks in this study involve vulnerability classes that are well-documented in public curricula and standard security education resources, including the MITRE CWE Top 25. Participants are not asked to develop working exploits or to interact with live systems; all code is generated in a local experimental environment. At the conclusion of each session, participants receive a structured debrief explaining the vulnerability class targeted by each task, the secure coding practice that addresses it, and the role of the multi-agent pipeline in identifying and mitigating the corresponding risk. This debrief serves both as a duty-of-care measure and as an educational benefit that partially compensates participants for their time.

%% ============================================================

\section{Methodology}

\subsection{Research Design}

This study employs a controlled between-subjects experiment to evaluate whether a structured multi-agent architecture produces measurably more secure code than a single-agent baseline when both systems operate on the same underlying model, tasks, and participants. The experimental design was chosen because it permits causal inference about the effect of the independent variable, architectural structure, on the dependent variables of vulnerability density, developer security reasoning quality, and task completion time. Observational and qualitative approaches were considered but rejected on the grounds that they cannot establish causality or support the statistical comparisons required to answer the research questions.

A between-subjects design was selected in preference to a within-subjects design for two reasons. First, participants who interact with the multi-agent pipeline are exposed to structured threat analysis and explicit vulnerability identification, which may cause lasting changes in security awareness that would contaminate any subsequent exposure to the baseline condition. Second, if the order were reversed and participants completed the baseline condition first, familiarity with the task format could artificially improve performance on the multi-agent condition. Either form of carryover would confound the comparison between conditions. Assigning each participant to a single condition eliminates this risk, at the cost of a larger required sample size.

The independent variable is the system architecture used to generate code. All other factors that might plausibly affect the security of the generated code are held constant: both conditions use GPT-4o as the underlying language model, both operate at a temperature of 0.2, and both receive the same four coding tasks. This design ensures that any statistically significant difference in security outcomes between conditions is attributable to architectural structure rather than to differences in model capability, generation stochasticity, or task difficulty.

\subsection{System Architecture}

\subsubsection*{Baseline System}

The baseline system implements a single-agent code generation pipeline using the OpenAI API directly, without any intermediate orchestration framework. When a participant submits a coding task, the system constructs a minimal prompt comprising a brief system message and the task description verbatim, and sends a single request to the GPT-4o API at temperature 0.2. The system prompt contains no security instructions, no structured review criteria, and no guidance toward defensive coding practices; this design choice is intentional and essential to the experimental validity of the baseline. The model response, containing the generated code and any accompanying explanation, is returned to the participant without modification or review.

Each baseline session is logged to a JSONL file with the following fields: \texttt{timestamp} (ISO 8601 UTC), \texttt{condition} (fixed value \texttt{"baseline"}), \texttt{participant\_id}, \texttt{model} (fixed value \texttt{"gpt-4o"}), \texttt{temperature} (fixed value \texttt{0.2}), \texttt{task}, \texttt{task\_id}, \texttt{task\_order}, \texttt{response}, and \texttt{duration\_seconds}. The \texttt{task\_order} field records the position of each task in the participant's randomised sequence and is included to support covariate analysis controlling for order effects at the analysis stage.

\subsubsection*{Multi-Agent System}

The multi-agent system implements a five-agent pipeline orchestrated using LangGraph, a directed acyclic graph execution framework for multi-step language model workflows. The five agents execute in a fixed sequence, with the output of each agent made available as structured input to subsequent agents. At defined points in the workflow, execution is interrupted and control is returned to the human participant, who reviews the agent's output and chooses whether to approve it, request a revision, or override it before the pipeline continues. This human-in-the-loop design is the principal architectural distinction between the multi-agent system and fully autonomous agentic pipelines.

The \textbf{Planner} receives the raw coding task and decomposes it into a structured implementation plan. Its output is a structured JSON object containing three fields: a list of implementation steps, a scope definition, and a set of security-relevant requirements implied by the task. The Planner does not access external knowledge sources; its function is to impose structure on the task before any code is generated.

The \textbf{Threat Modeller} receives the Planner output and identifies security threats relevant to the task before any code is written. It queries a local ChromaDB vector database containing the MITRE CWE Top 25 Most Dangerous Software Weaknesses \citep{mitre2024} using an Advanced Retrieval-Augmented Generation pipeline. This pipeline proceeds in four stages: first, a query rewrite step in which GPT-4o expands the task description into security-relevant search terms; second, a metadata filter that restricts the vector search to CWE entries relevant to the task domain; third, a similarity search over ChromaDB returning the top ten most relevant CWE entries; and fourth, an LLM re-ranking step in which GPT-4o selects the three most relevant entries for injection into the agent prompt. In parallel, the Threat Modeller queries the NIST National Vulnerability Database REST API (v2) via a Model Context Protocol server to retrieve recent Common Vulnerabilities and Exposures records associated with keywords derived from the task. The Threat Modeller's output is a structured JSON list of threats, each comprising a CWE identifier, a severity rating, and a recommended mitigation strategy.

The \textbf{Code Generator} receives both the Planner and Threat Modeller outputs and produces the implementation code informed by both. It does not access external knowledge sources; the security context provided by the Threat Modeller's structured output is the mechanism by which security reasoning is embedded in code generation. The Code Generator's output is a structured JSON object containing a code block and a natural-language explanation of the security decisions made during implementation.

The \textbf{Code Reviewer} receives the generated code and conducts an independent security review in two stages. First, it invokes a Bandit static analysis MCP server, which executes the \texttt{bandit -f json} command on the generated code and returns structured JSON findings including issue severity, confidence, and CWE mapping. Second, it uses these findings as structured input alongside the code to produce its own assessment. The Code Reviewer's output is a structured JSON list of security issues, each with a CWE identifier, severity rating, and suggested fix. The Code Reviewer operates independently of the Verifier; it does not share its Bandit invocation or findings with the subsequent agent.

The \textbf{Verifier} performs a final, independent validation of the generated code against the original threat model. It independently queries the same ChromaDB CWE corpus using the same Advanced RAG pipeline as the Threat Modeller, arriving at its own assessment without access to the Threat Modeller's findings. It independently invokes the Bandit MCP server on the generated code. It also executes the code in a sandboxed subprocess via a third MCP server, enforcing a strict timeout and prohibiting network access and file system writes, to verify that the code runs without runtime errors. The Verifier's output is a structured JSON object containing a pass or fail verdict for each threat identified in the original threat model, the static analysis findings from its independent Bandit run, and the execution result from the sandbox.

Session records in the multi-agent condition extend the baseline log schema with per-agent fields recording each agent's structured output and the participant's decision at each human-in-the-loop interrupt. The \texttt{task\_id} and \texttt{task\_order} fields are present in both conditions to support cross-condition covariate analysis.

\subsection{Participants}

The minimum required sample size was determined by an a priori power analysis conducted in G*Power 3.1, specifying 80 per cent power, a two-tailed alpha of 0.05, and a medium effect size of Cohen's $d = 0.5$. The effect size of $d = 0.5$ was adopted from \citet{perry2023}, whose controlled study of AI-assisted code generation provides the closest existing empirical precedent for the effect magnitudes expected in the present study. The power analysis returned a minimum total $n$ of 52, requiring 26 participants per condition.

Participants are recruited from computer science and software engineering students enrolled at the University of Exeter. Inclusion criteria are age 18 or over, basic working proficiency in Python, and current enrolment at the University of Exeter. Participants with no Python experience are excluded on the grounds that they would be unable to engage meaningfully with the coding tasks or to evaluate the quality of AI-generated code during the human oversight steps. Recruitment proceeds via departmental email lists and course notice boards. Participants are not informed of the specific research hypothesis prior to the session to prevent demand characteristics from influencing their behaviour, consistent with standard practice in within-session experimental designs.

Ethics approval was obtained from the University of Exeter Engineering, Mathematics and Physical Sciences Research Ethics Committee (Application ID: 13362509) prior to any recruitment activity.

\subsection{Tasks}

The experiment uses four Python coding tasks, each targeting a distinct vulnerability category from the MITRE CWE Top 25 Most Dangerous Software Weaknesses \citep{mitre2024}. The tasks are as follows. The first task requires participants to write a Python function that accepts a username and password, checks them against credentials stored in a SQLite database, and returns a Boolean login result; this task targets CWE-916 (use of a password hash with insufficient computational effort) and CWE-312 (cleartext storage of sensitive information). The second task requires a function that accepts a filename from the user and reads and returns the file contents; this task targets CWE-22 (path traversal). The third task requires a function that connects to a SQLite database, accepts a username as input, and returns the corresponding record from a \texttt{users} table; this task targets CWE-89 (SQL injection). The fourth task requires a function that accepts a string from a web form and displays it back to the user on a webpage; this task targets CWE-79 (cross-site scripting).

Task descriptions are deliberately neutral with respect to security. No task description mentions vulnerabilities, secure coding practices, or defensive programming, so that neither condition receives a security cue from the task text itself. The contrast between conditions arises entirely from the pipeline through which the task is processed.

Task order is randomised independently for each participant using a random permutation of the four task identifiers, and the assigned order is recorded in every session log. Randomisation controls for learning effects, in which awareness of vulnerability patterns elicited by one task might influence participant behaviour on subsequent tasks. Each task is allocated a maximum of 15 minutes of session time.

\subsection{Evaluation Metrics and Scoring}

\paragraph{Security.} The primary outcome measure is vulnerability density, defined as the number of CWE-mapped vulnerabilities identified in the generated code per 100 lines of code. Vulnerabilities are identified using Bandit static analysis, the established Python security linter maintained by the Python Security project, which produces structured JSON output with issue severity, confidence rating, and CWE identifier. Bandit was selected for three reasons: it is the standard tool for Python security static analysis; its output is structured and CWE-mapped, enabling consistent comparison across conditions; and it is the same tool used within the multi-agent pipeline itself, ensuring that the evaluation metric is consistent with the internal security checking performed during code generation. All Bandit assessments are conducted on the final code output produced by each condition, using a fixed configuration to ensure comparability.

\paragraph{Reasoning quality.} Developer security reasoning is assessed using a structured rubric applied to the participant's written justifications and decisions produced during the session. Each response is scored on a scale of 0 to 3 across three criteria: accuracy of threat identification, quality of design justification, and evidence of critical evaluation of AI-generated recommendations. Rubric application is conducted by two independent raters, with inter-rater reliability quantified using Cohen's kappa. A target kappa of 0.7 or above, representing substantial agreement, is required before scores are used in analysis. Where the kappa falls below this threshold, a reconciliation procedure is conducted before analysis proceeds.

\paragraph{Productivity.} Task completion time, measured in minutes and logged automatically by the session management system, is used as the productivity metric. The study's success criterion specifies that the multi-agent condition should not incur a productivity cost exceeding 20 per cent relative to the baseline.

\subsection{Statistical Analysis}

The Mann-Whitney U test is used as the primary inferential test for the vulnerability density outcome, because vulnerability counts per task are expected to be non-normally distributed and bounded below by zero, making parametric tests inappropriate. The Wilcoxon signed-rank test is applied to reasoning quality scores where appropriate. Effect sizes are reported alongside all inferential test results: Cohen's $d$ is reported for comparisons involving continuous measures, and the rank-biserial correlation $r$ is reported for Mann-Whitney U results. All statistical analysis is conducted in Python using the \texttt{scipy.stats} module. The alpha level is set to 0.05 throughout, with no adjustment for multiple comparisons given that the primary and secondary outcomes are specified a priori and are conceptually distinct.

\subsection{Limitations and Mitigations}

Several limitations of the present methodology warrant explicit acknowledgement. First, the participant sample consists of computer science and software engineering students at a single institution, rather than professional developers. Students may differ from industry practitioners in their Python proficiency, their familiarity with security-relevant coding patterns, and their propensity to engage critically with AI-generated suggestions. This limits the external validity of the findings, and the degree to which conclusions can be generalised to professional software development contexts is discussed in the relevant section. The homogeneity of the sample, a consequence of single-institution recruitment, further constrains the representativeness of the findings.

Second, the reasoning quality rubric, despite its structured criteria and inter-rater reliability checks, requires evaluative judgement that introduces a degree of subjectivity not present in the automated vulnerability density metric. Discrepancies between raters are resolved by a defined reconciliation procedure, but this cannot eliminate the possibility that rater expectations about the multi-agent condition influence their scoring.

Third, the four tasks used in this study are short, self-contained functions targeting single vulnerability classes. \citet{shukla2025} find that iterative AI assistance on larger, more complex code artefacts can produce compounding security degradation that would not be apparent in single-function tasks. The ecological validity of short tasks relative to real-world software development is therefore limited; whether the findings extend to more complex multi-function systems cannot be established from this study alone.

Fourth, the quality of human oversight in the multi-agent condition is not standardised across participants. \citet{kennedy2025} report that humans in AI oversight roles provide correct oversight only approximately half the time under typical conditions, with failures driven by pressure to approve rather than scrutinise. Participants who habitually approve agent outputs without genuine engagement would artificially inflate reasoning scores and might reduce the apparent benefit of the human-in-the-loop design. To mitigate this, participants complete a calibration task at the start of the multi-agent session that requires them to identify a deliberate error introduced into an agent output, enabling passive acceptance patterns to be detected and flagged in the analysis.

\subsection{Ethical Considerations}

Ethics approval was obtained from the University of Exeter Engineering, Mathematics and Physical Sciences Research Ethics Committee (Application ID: 13362509) prior to any recruitment activity. No participant data was collected before approval was confirmed.

Participation is entirely voluntary. Prior to each session, participants receive a participant information sheet detailing the purpose of the study, the nature of the tasks, the data collected, and their rights. Electronic informed consent is obtained before the session begins. Participants are made aware that they may withdraw from the study and request deletion of their data within two weeks of their session by quoting their assigned participant identifier, without giving reasons and without penalty.

Data are anonymised at the point of collection. Each participant is assigned a unique identifier (P01, P02, and so on) before the session begins, and no names, email addresses, or other personally identifying information are stored in session logs. All session data are stored on University of Exeter OneDrive storage held by the project supervisor, accessible only to the researcher and supervisor. Session log files are excluded from version control via the project \texttt{.gitignore} configuration.

The coding tasks in this study involve vulnerability classes that are well-documented in public curricula and standard security education resources, including the MITRE CWE Top 25. Participants are not asked to develop working exploits or to interact with live systems; all code is generated in a local experimental environment. At the conclusion of each session, participants receive a structured debrief explaining the vulnerability class targeted by each task, the secure coding practice that addresses it, and the role of the multi-agent pipeline in identifying and mitigating the corresponding risk. This debrief serves both as a duty-of-care measure and as an educational benefit that partially compensates participants for their time.

%% ============================================================

\section{Methodology}

\subsection{Research Design}

This study employs a controlled between-subjects experiment to evaluate whether a structured multi-agent architecture produces measurably more secure code than a single-agent baseline when both systems operate on the same underlying model, tasks, and participants. The experimental design was chosen because it permits causal inference about the effect of the independent variable, architectural structure, on the dependent variables of vulnerability density, developer security reasoning quality, and task completion time. Observational and qualitative approaches were considered but rejected on the grounds that they cannot establish causality or support the statistical comparisons required to answer the research questions.

A between-subjects design was selected in preference to a within-subjects design for two reasons. First, participants who interact with the multi-agent pipeline are exposed to structured threat analysis and explicit vulnerability identification, which may cause lasting changes in security awareness that would contaminate any subsequent exposure to the baseline condition. Second, if the order were reversed and participants completed the baseline condition first, familiarity with the task format could artificially improve performance on the multi-agent condition. Either form of carryover would confound the comparison between conditions. Assigning each participant to a single condition eliminates this risk, at the cost of a larger required sample size.

The independent variable is the system architecture used to generate code. All other factors that might plausibly affect the security of the generated code are held constant: both conditions use GPT-4o as the underlying language model, both operate at a temperature of 0.2, and both receive the same four coding tasks. This design ensures that any statistically significant difference in security outcomes between conditions is attributable to architectural structure rather than to differences in model capability, generation stochasticity, or task difficulty.

\subsection{System Architecture}

\subsubsection*{Baseline System}

The baseline system implements a single-agent code generation pipeline using the OpenAI API directly, without any intermediate orchestration framework. When a participant submits a coding task, the system constructs a minimal prompt comprising a brief system message and the task description verbatim, and sends a single request to the GPT-4o API at temperature 0.2. The system prompt contains no security instructions, no structured review criteria, and no guidance toward defensive coding practices; this design choice is intentional and essential to the experimental validity of the baseline. The model response, containing the generated code and any accompanying explanation, is returned to the participant without modification or review.

Each baseline session is logged to a JSONL file with the following fields: \texttt{timestamp} (ISO 8601 UTC), \texttt{condition} (fixed value \texttt{"baseline"}), \texttt{participant\_id}, \texttt{model} (fixed value \texttt{"gpt-4o"}), \texttt{temperature} (fixed value \texttt{0.2}), \texttt{task}, \texttt{task\_id}, \texttt{task\_order}, \texttt{response}, and \texttt{duration\_seconds}. The \texttt{task\_order} field records the position of each task in the participant's randomised sequence and is included to support covariate analysis controlling for order effects at the analysis stage.

\subsubsection*{Multi-Agent System}

The multi-agent system implements a five-agent pipeline orchestrated using LangGraph, a directed acyclic graph execution framework for multi-step language model workflows. The five agents execute in a fixed sequence, with the output of each agent made available as structured input to subsequent agents. At defined points in the workflow, execution is interrupted and control is returned to the human participant, who reviews the agent's output and chooses whether to approve it, request a revision, or override it before the pipeline continues. This human-in-the-loop design is the principal architectural distinction between the multi-agent system and fully autonomous agentic pipelines.

The \textbf{Planner} receives the raw coding task and decomposes it into a structured implementation plan. Its output is a structured JSON object containing three fields: a list of implementation steps, a scope definition, and a set of security-relevant requirements implied by the task. The Planner does not access external knowledge sources; its function is to impose structure on the task before any code is generated.

The \textbf{Threat Modeller} receives the Planner output and identifies security threats relevant to the task before any code is written. It queries a local ChromaDB vector database containing the MITRE CWE Top 25 Most Dangerous Software Weaknesses \citep{mitre2024} using an Advanced Retrieval-Augmented Generation pipeline. This pipeline proceeds in four stages: first, a query rewrite step in which GPT-4o expands the task description into security-relevant search terms; second, a metadata filter that restricts the vector search to CWE entries relevant to the task domain; third, a similarity search over ChromaDB returning the top ten most relevant CWE entries; and fourth, an LLM re-ranking step in which GPT-4o selects the three most relevant entries for injection into the agent prompt. In parallel, the Threat Modeller queries the NIST National Vulnerability Database REST API (v2) via a Model Context Protocol server to retrieve recent Common Vulnerabilities and Exposures records associated with keywords derived from the task. The Threat Modeller's output is a structured JSON list of threats, each comprising a CWE identifier, a severity rating, and a recommended mitigation strategy.

The \textbf{Code Generator} receives both the Planner and Threat Modeller outputs and produces the implementation code informed by both. It does not access external knowledge sources; the security context provided by the Threat Modeller's structured output is the mechanism by which security reasoning is embedded in code generation. The Code Generator's output is a structured JSON object containing a code block and a natural-language explanation of the security decisions made during implementation.

The \textbf{Code Reviewer} receives the generated code and conducts an independent security review in two stages. First, it invokes a Bandit static analysis MCP server, which executes the \texttt{bandit -f json} command on the generated code and returns structured JSON findings including issue severity, confidence, and CWE mapping. Second, it uses these findings as structured input alongside the code to produce its own assessment. The Code Reviewer's output is a structured JSON list of security issues, each with a CWE identifier, severity rating, and suggested fix. The Code Reviewer operates independently of the Verifier; it does not share its Bandit invocation or findings with the subsequent agent.

The \textbf{Verifier} performs a final, independent validation of the generated code against the original threat model. It independently queries the same ChromaDB CWE corpus using the same Advanced RAG pipeline as the Threat Modeller, arriving at its own assessment without access to the Threat Modeller's findings. It independently invokes the Bandit MCP server on the generated code. It also executes the code in a sandboxed subprocess via a third MCP server, enforcing a strict timeout and prohibiting network access and file system writes, to verify that the code runs without runtime errors. The Verifier's output is a structured JSON object containing a pass or fail verdict for each threat identified in the original threat model, the static analysis findings from its independent Bandit run, and the execution result from the sandbox.

Session records in the multi-agent condition extend the baseline log schema with per-agent fields recording each agent's structured output and the participant's decision at each human-in-the-loop interrupt. The \texttt{task\_id} and \texttt{task\_order} fields are present in both conditions to support cross-condition covariate analysis.

\subsection{Participants}

The minimum required sample size was determined by an a priori power analysis conducted in G*Power 3.1, specifying 80 per cent power, a two-tailed alpha of 0.05, and a medium effect size of Cohen's $d = 0.5$. The effect size of $d = 0.5$ was adopted from \citet{perry2023}, whose controlled study of AI-assisted code generation provides the closest existing empirical precedent for the effect magnitudes expected in the present study. The power analysis returned a minimum total $n$ of 52, requiring 26 participants per condition.

Participants are recruited from computer science and software engineering students enrolled at the University of Exeter. Inclusion criteria are age 18 or over, basic working proficiency in Python, and current enrolment at the University of Exeter. Participants with no Python experience are excluded on the grounds that they would be unable to engage meaningfully with the coding tasks or to evaluate the quality of AI-generated code during the human oversight steps. Recruitment proceeds via departmental email lists and course notice boards. Participants are not informed of the specific research hypothesis prior to the session to prevent demand characteristics from influencing their behaviour, consistent with standard practice in within-session experimental designs.

Ethics approval was obtained from the University of Exeter Engineering, Mathematics and Physical Sciences Research Ethics Committee (Application ID: 13362509) prior to any recruitment activity.

\subsection{Tasks}

The experiment uses four Python coding tasks, each targeting a distinct vulnerability category from the MITRE CWE Top 25 Most Dangerous Software Weaknesses \citep{mitre2024}. The tasks are as follows. The first task requires participants to write a Python function that accepts a username and password, checks them against credentials stored in a SQLite database, and returns a Boolean login result; this task targets CWE-916 (use of a password hash with insufficient computational effort) and CWE-312 (cleartext storage of sensitive information). The second task requires a function that accepts a filename from the user and reads and returns the file contents; this task targets CWE-22 (path traversal). The third task requires a function that connects to a SQLite database, accepts a username as input, and returns the corresponding record from a \texttt{users} table; this task targets CWE-89 (SQL injection). The fourth task requires a function that accepts a string from a web form and displays it back to the user on a webpage; this task targets CWE-79 (cross-site scripting).

Task descriptions are deliberately neutral with respect to security. No task description mentions vulnerabilities, secure coding practices, or defensive programming, so that neither condition receives a security cue from the task text itself. The contrast between conditions arises entirely from the pipeline through which the task is processed.

Task order is randomised independently for each participant using a random permutation of the four task identifiers, and the assigned order is recorded in every session log. Randomisation controls for learning effects, in which awareness of vulnerability patterns elicited by one task might influence participant behaviour on subsequent tasks. Each task is allocated a maximum of 15 minutes of session time.

\subsection{Evaluation Metrics and Scoring}

\paragraph{Security.} The primary outcome measure is vulnerability density, defined as the number of CWE-mapped vulnerabilities identified in the generated code per 100 lines of code. Vulnerabilities are identified using Bandit static analysis, the established Python security linter maintained by the Python Security project, which produces structured JSON output with issue severity, confidence rating, and CWE identifier. Bandit was selected for three reasons: it is the standard tool for Python security static analysis; its output is structured and CWE-mapped, enabling consistent comparison across conditions; and it is the same tool used within the multi-agent pipeline itself, ensuring that the evaluation metric is consistent with the internal security checking performed during code generation. All Bandit assessments are conducted on the final code output produced by each condition, using a fixed configuration to ensure comparability.

\paragraph{Reasoning quality.} Developer security reasoning is assessed using a structured rubric applied to the participant's written justifications and decisions produced during the session. Each response is scored on a scale of 0 to 3 across three criteria: accuracy of threat identification, quality of design justification, and evidence of critical evaluation of AI-generated recommendations. Rubric application is conducted by two independent raters, with inter-rater reliability quantified using Cohen's kappa. A target kappa of 0.7 or above, representing substantial agreement, is required before scores are used in analysis. Where the kappa falls below this threshold, a reconciliation procedure is conducted before analysis proceeds.

\paragraph{Productivity.} Task completion time, measured in minutes and logged automatically by the session management system, is used as the productivity metric. The study's success criterion specifies that the multi-agent condition should not incur a productivity cost exceeding 20 per cent relative to the baseline.

\subsection{Statistical Analysis}

The Mann-Whitney U test is used as the primary inferential test for the vulnerability density outcome, because vulnerability counts per task are expected to be non-normally distributed and bounded below by zero, making parametric tests inappropriate. The Wilcoxon signed-rank test is applied to reasoning quality scores where appropriate. Effect sizes are reported alongside all inferential test results: Cohen's $d$ is reported for comparisons involving continuous measures, and the rank-biserial correlation $r$ is reported for Mann-Whitney U results. All statistical analysis is conducted in Python using the \texttt{scipy.stats} module. The alpha level is set to 0.05 throughout, with no adjustment for multiple comparisons given that the primary and secondary outcomes are specified a priori and are conceptually distinct.

\subsection{Limitations and Mitigations}

Several limitations of the present methodology warrant explicit acknowledgement. First, the participant sample consists of computer science and software engineering students at a single institution, rather than professional developers. Students may differ from industry practitioners in their Python proficiency, their familiarity with security-relevant coding patterns, and their propensity to engage critically with AI-generated suggestions. This limits the external validity of the findings, and the degree to which conclusions can be generalised to professional software development contexts is discussed in the relevant section. The homogeneity of the sample, a consequence of single-institution recruitment, further constrains the representativeness of the findings.

Second, the reasoning quality rubric, despite its structured criteria and inter-rater reliability checks, requires evaluative judgement that introduces a degree of subjectivity not present in the automated vulnerability density metric. Discrepancies between raters are resolved by a defined reconciliation procedure, but this cannot eliminate the possibility that rater expectations about the multi-agent condition influence their scoring.

Third, the four tasks used in this study are short, self-contained functions targeting single vulnerability classes. \citet{shukla2025} find that iterative AI assistance on larger, more complex code artefacts can produce compounding security degradation that would not be apparent in single-function tasks. The ecological validity of short tasks relative to real-world software development is therefore limited; whether the findings extend to more complex multi-function systems cannot be established from this study alone.

Fourth, the quality of human oversight in the multi-agent condition is not standardised across participants. \citet{kennedy2025} report that humans in AI oversight roles provide correct oversight only approximately half the time under typical conditions, with failures driven by pressure to approve rather than scrutinise. Participants who habitually approve agent outputs without genuine engagement would artificially inflate reasoning scores and might reduce the apparent benefit of the human-in-the-loop design. To mitigate this, participants complete a calibration task at the start of the multi-agent session that requires them to identify a deliberate error introduced into an agent output, enabling passive acceptance patterns to be detected and flagged in the analysis.

\subsection{Ethical Considerations}

Ethics approval was obtained from the University of Exeter Engineering, Mathematics and Physical Sciences Research Ethics Committee (Application ID: 13362509) prior to any recruitment activity. No participant data was collected before approval was confirmed.

Participation is entirely voluntary. Prior to each session, participants receive a participant information sheet detailing the purpose of the study, the nature of the tasks, the data collected, and their rights. Electronic informed consent is obtained before the session begins. Participants are made aware that they may withdraw from the study and request deletion of their data within two weeks of their session by quoting their assigned participant identifier, without giving reasons and without penalty.

Data are anonymised at the point of collection. Each participant is assigned a unique identifier (P01, P02, and so on) before the session begins, and no names, email addresses, or other personally identifying information are stored in session logs. All session data are stored on University of Exeter OneDrive storage held by the project supervisor, accessible only to the researcher and supervisor. Session log files are excluded from version control via the project \texttt{.gitignore} configuration.

The coding tasks in this study involve vulnerability classes that are well-documented in public curricula and standard security education resources, including the MITRE CWE Top 25. Participants are not asked to develop working exploits or to interact with live systems; all code is generated in a local experimental environment. At the conclusion of each session, participants receive a structured debrief explaining the vulnerability class targeted by each task, the secure coding practice that addresses it, and the role of the multi-agent pipeline in identifying and mitigating the corresponding risk. This debrief serves both as a duty-of-care measure and as an educational benefit that partially compensates participants for their time.

%% ============================================================

\section{Methodology}

\subsection{Research Design}

This study employs a controlled between-subjects experiment to evaluate whether a structured multi-agent architecture produces measurably more secure code than a single-agent baseline when both systems operate on the same underlying model, tasks, and participants. The experimental design was chosen because it permits causal inference about the effect of the independent variable, architectural structure, on the dependent variables of vulnerability density, developer security reasoning quality, and task completion time. Observational and qualitative approaches were considered but rejected on the grounds that they cannot establish causality or support the statistical comparisons required to answer the research questions.

A between-subjects design was selected in preference to a within-subjects design for two reasons. First, participants who interact with the multi-agent pipeline are exposed to structured threat analysis and explicit vulnerability identification, which may cause lasting changes in security awareness that would contaminate any subsequent exposure to the baseline condition. Second, if the order were reversed and participants completed the baseline condition first, familiarity with the task format could artificially improve performance on the multi-agent condition. Either form of carryover would confound the comparison between conditions. Assigning each participant to a single condition eliminates this risk, at the cost of a larger required sample size.

The independent variable is the system architecture used to generate code. All other factors that might plausibly affect the security of the generated code are held constant: both conditions use GPT-4o as the underlying language model, both operate at a temperature of 0.2, and both receive the same four coding tasks. This design ensures that any statistically significant difference in security outcomes between conditions is attributable to architectural structure rather than to differences in model capability, generation stochasticity, or task difficulty.

\subsection{System Architecture}

\subsubsection*{Baseline System}

The baseline system implements a single-agent code generation pipeline using the OpenAI API directly, without any intermediate orchestration framework. When a participant submits a coding task, the system constructs a minimal prompt comprising a brief system message and the task description verbatim, and sends a single request to the GPT-4o API at temperature 0.2. The system prompt contains no security instructions, no structured review criteria, and no guidance toward defensive coding practices; this design choice is intentional and essential to the experimental validity of the baseline. The model response, containing the generated code and any accompanying explanation, is returned to the participant without modification or review.

Each baseline session is logged to a JSONL file with the following fields: \texttt{timestamp} (ISO 8601 UTC), \texttt{condition} (fixed value \texttt{"baseline"}), \texttt{participant\_id}, \texttt{model} (fixed value \texttt{"gpt-4o"}), \texttt{temperature} (fixed value \texttt{0.2}), \texttt{task}, \texttt{task\_id}, \texttt{task\_order}, \texttt{response}, and \texttt{duration\_seconds}. The \texttt{task\_order} field records the position of each task in the participant's randomised sequence and is included to support covariate analysis controlling for order effects at the analysis stage.

\subsubsection*{Multi-Agent System}

The multi-agent system implements a five-agent pipeline orchestrated using LangGraph, a directed acyclic graph execution framework for multi-step language model workflows. The five agents execute in a fixed sequence, with the output of each agent made available as structured input to subsequent agents. At defined points in the workflow, execution is interrupted and control is returned to the human participant, who reviews the agent's output and chooses whether to approve it, request a revision, or override it before the pipeline continues. This human-in-the-loop design is the principal architectural distinction between the multi-agent system and fully autonomous agentic pipelines.

The \textbf{Planner} receives the raw coding task and decomposes it into a structured implementation plan. Its output is a structured JSON object containing three fields: a list of implementation steps, a scope definition, and a set of security-relevant requirements implied by the task. The Planner does not access external knowledge sources; its function is to impose structure on the task before any code is generated.

The \textbf{Threat Modeller} receives the Planner output and identifies security threats relevant to the task before any code is written. It queries a local ChromaDB vector database containing the MITRE CWE Top 25 Most Dangerous Software Weaknesses \citep{mitre2024} using an Advanced Retrieval-Augmented Generation pipeline. This pipeline proceeds in four stages: first, a query rewrite step in which GPT-4o expands the task description into security-relevant search terms; second, a metadata filter that restricts the vector search to CWE entries relevant to the task domain; third, a similarity search over ChromaDB returning the top ten most relevant CWE entries; and fourth, an LLM re-ranking step in which GPT-4o selects the three most relevant entries for injection into the agent prompt. In parallel, the Threat Modeller queries the NIST National Vulnerability Database REST API (v2) via a Model Context Protocol server to retrieve recent Common Vulnerabilities and Exposures records associated with keywords derived from the task. The Threat Modeller's output is a structured JSON list of threats, each comprising a CWE identifier, a severity rating, and a recommended mitigation strategy.

The \textbf{Code Generator} receives both the Planner and Threat Modeller outputs and produces the implementation code informed by both. It does not access external knowledge sources; the security context provided by the Threat Modeller's structured output is the mechanism by which security reasoning is embedded in code generation. The Code Generator's output is a structured JSON object containing a code block and a natural-language explanation of the security decisions made during implementation.

The \textbf{Code Reviewer} receives the generated code and conducts an independent security review in two stages. First, it invokes a Bandit static analysis MCP server, which executes the \texttt{bandit -f json} command on the generated code and returns structured JSON findings including issue severity, confidence, and CWE mapping. Second, it uses these findings as structured input alongside the code to produce its own assessment. The Code Reviewer's output is a structured JSON list of security issues, each with a CWE identifier, severity rating, and suggested fix. The Code Reviewer operates independently of the Verifier; it does not share its Bandit invocation or findings with the subsequent agent.

The \textbf{Verifier} performs a final, independent validation of the generated code against the original threat model. It independently queries the same ChromaDB CWE corpus using the same Advanced RAG pipeline as the Threat Modeller, arriving at its own assessment without access to the Threat Modeller's findings. It independently invokes the Bandit MCP server on the generated code. It also executes the code in a sandboxed subprocess via a third MCP server, enforcing a strict timeout and prohibiting network access and file system writes, to verify that the code runs without runtime errors. The Verifier's output is a structured JSON object containing a pass or fail verdict for each threat identified in the original threat model, the static analysis findings from its independent Bandit run, and the execution result from the sandbox.

Session records in the multi-agent condition extend the baseline log schema with per-agent fields recording each agent's structured output and the participant's decision at each human-in-the-loop interrupt. The \texttt{task\_id} and \texttt{task\_order} fields are present in both conditions to support cross-condition covariate analysis.

\subsection{Participants}

The minimum required sample size was determined by an a priori power analysis conducted in G*Power 3.1, specifying 80 per cent power, a two-tailed alpha of 0.05, and a medium effect size of Cohen's $d = 0.5$. The effect size of $d = 0.5$ was adopted from \citet{perry2023}, whose controlled study of AI-assisted code generation provides the closest existing empirical precedent for the effect magnitudes expected in the present study. The power analysis returned a minimum total $n$ of 52, requiring 26 participants per condition.

Participants are recruited from computer science and software engineering students enrolled at the University of Exeter. Inclusion criteria are age 18 or over, basic working proficiency in Python, and current enrolment at the University of Exeter. Participants with no Python experience are excluded on the grounds that they would be unable to engage meaningfully with the coding tasks or to evaluate the quality of AI-generated code during the human oversight steps. Recruitment proceeds via departmental email lists and course notice boards. Participants are not informed of the specific research hypothesis prior to the session to prevent demand characteristics from influencing their behaviour, consistent with standard practice in within-session experimental designs.

Ethics approval was obtained from the University of Exeter Engineering, Mathematics and Physical Sciences Research Ethics Committee (Application ID: 13362509) prior to any recruitment activity.

\subsection{Tasks}

The experiment uses four Python coding tasks, each targeting a distinct vulnerability category from the MITRE CWE Top 25 Most Dangerous Software Weaknesses \citep{mitre2024}. The tasks are as follows. The first task requires participants to write a Python function that accepts a username and password, checks them against credentials stored in a SQLite database, and returns a Boolean login result; this task targets CWE-916 (use of a password hash with insufficient computational effort) and CWE-312 (cleartext storage of sensitive information). The second task requires a function that accepts a filename from the user and reads and returns the file contents; this task targets CWE-22 (path traversal). The third task requires a function that connects to a SQLite database, accepts a username as input, and returns the corresponding record from a \texttt{users} table; this task targets CWE-89 (SQL injection). The fourth task requires a function that accepts a string from a web form and displays it back to the user on a webpage; this task targets CWE-79 (cross-site scripting).

Task descriptions are deliberately neutral with respect to security. No task description mentions vulnerabilities, secure coding practices, or defensive programming, so that neither condition receives a security cue from the task text itself. The contrast between conditions arises entirely from the pipeline through which the task is processed.

Task order is randomised independently for each participant using a random permutation of the four task identifiers, and the assigned order is recorded in every session log. Randomisation controls for learning effects, in which awareness of vulnerability patterns elicited by one task might influence participant behaviour on subsequent tasks. Each task is allocated a maximum of 15 minutes of session time.

\subsection{Evaluation Metrics and Scoring}

\paragraph{Security.} The primary outcome measure is vulnerability density, defined as the number of CWE-mapped vulnerabilities identified in the generated code per 100 lines of code. Vulnerabilities are identified using Bandit static analysis, the established Python security linter maintained by the Python Security project, which produces structured JSON output with issue severity, confidence rating, and CWE identifier. Bandit was selected for three reasons: it is the standard tool for Python security static analysis; its output is structured and CWE-mapped, enabling consistent comparison across conditions; and it is the same tool used within the multi-agent pipeline itself, ensuring that the evaluation metric is consistent with the internal security checking performed during code generation. All Bandit assessments are conducted on the final code output produced by each condition, using a fixed configuration to ensure comparability.

\paragraph{Reasoning quality.} Developer security reasoning is assessed using a structured rubric applied to the participant's written justifications and decisions produced during the session. Each response is scored on a scale of 0 to 3 across three criteria: accuracy of threat identification, quality of design justification, and evidence of critical evaluation of AI-generated recommendations. Rubric application is conducted by two independent raters, with inter-rater reliability quantified using Cohen's kappa. A target kappa of 0.7 or above, representing substantial agreement, is required before scores are used in analysis. Where the kappa falls below this threshold, a reconciliation procedure is conducted before analysis proceeds.

\paragraph{Productivity.} Task completion time, measured in minutes and logged automatically by the session management system, is used as the productivity metric. The study's success criterion specifies that the multi-agent condition should not incur a productivity cost exceeding 20 per cent relative to the baseline.

\subsection{Statistical Analysis}

The Mann-Whitney U test is used as the primary inferential test for the vulnerability density outcome, because vulnerability counts per task are expected to be non-normally distributed and bounded below by zero, making parametric tests inappropriate. The Wilcoxon signed-rank test is applied to reasoning quality scores where appropriate. Effect sizes are reported alongside all inferential test results: Cohen's $d$ is reported for comparisons involving continuous measures, and the rank-biserial correlation $r$ is reported for Mann-Whitney U results. All statistical analysis is conducted in Python using the \texttt{scipy.stats} module. The alpha level is set to 0.05 throughout, with no adjustment for multiple comparisons given that the primary and secondary outcomes are specified a priori and are conceptually distinct.

\subsection{Limitations and Mitigations}

Several limitations of the present methodology warrant explicit acknowledgement. First, the participant sample consists of computer science and software engineering students at a single institution, rather than professional developers. Students may differ from industry practitioners in their Python proficiency, their familiarity with security-relevant coding patterns, and their propensity to engage critically with AI-generated suggestions. This limits the external validity of the findings, and the degree to which conclusions can be generalised to professional software development contexts is discussed in the relevant section. The homogeneity of the sample, a consequence of single-institution recruitment, further constrains the representativeness of the findings.

Second, the reasoning quality rubric, despite its structured criteria and inter-rater reliability checks, requires evaluative judgement that introduces a degree of subjectivity not present in the automated vulnerability density metric. Discrepancies between raters are resolved by a defined reconciliation procedure, but this cannot eliminate the possibility that rater expectations about the multi-agent condition influence their scoring.

Third, the four tasks used in this study are short, self-contained functions targeting single vulnerability classes. \citet{shukla2025} find that iterative AI assistance on larger, more complex code artefacts can produce compounding security degradation that would not be apparent in single-function tasks. The ecological validity of short tasks relative to real-world software development is therefore limited; whether the findings extend to more complex multi-function systems cannot be established from this study alone.

Fourth, the quality of human oversight in the multi-agent condition is not standardised across participants. \citet{kennedy2025} report that humans in AI oversight roles provide correct oversight only approximately half the time under typical conditions, with failures driven by pressure to approve rather than scrutinise. Participants who habitually approve agent outputs without genuine engagement would artificially inflate reasoning scores and might reduce the apparent benefit of the human-in-the-loop design. To mitigate this, participants complete a calibration task at the start of the multi-agent session that requires them to identify a deliberate error introduced into an agent output, enabling passive acceptance patterns to be detected and flagged in the analysis.

\subsection{Ethical Considerations}

Ethics approval was obtained from the University of Exeter Engineering, Mathematics and Physical Sciences Research Ethics Committee (Application ID: 13362509) prior to any recruitment activity. No participant data was collected before approval was confirmed.

Participation is entirely voluntary. Prior to each session, participants receive a participant information sheet detailing the purpose of the study, the nature of the tasks, the data collected, and their rights. Electronic informed consent is obtained before the session begins. Participants are made aware that they may withdraw from the study and request deletion of their data within two weeks of their session by quoting their assigned participant identifier, without giving reasons and without penalty.

Data are anonymised at the point of collection. Each participant is assigned a unique identifier (P01, P02, and so on) before the session begins, and no names, email addresses, or other personally identifying information are stored in session logs. All session data are stored on University of Exeter OneDrive storage held by the project supervisor, accessible only to the researcher and supervisor. Session log files are excluded from version control via the project \texttt{.gitignore} configuration.

The coding tasks in this study involve vulnerability classes that are well-documented in public curricula and standard security education resources, including the MITRE CWE Top 25. Participants are not asked to develop working exploits or to interact with live systems; all code is generated in a local experimental environment. At the conclusion of each session, participants receive a structured debrief explaining the vulnerability class targeted by each task, the secure coding practice that addresses it, and the role of the multi-agent pipeline in identifying and mitigating the corresponding risk. This debrief serves both as a duty-of-care measure and as an educational benefit that partially compensates participants for their time.
